\newcommand{\AstNode}{\mathord{\vcenter{\hbox{\begin{tikzpicture}[x=1ex,y=1ex,line width=.06ex]
\draw (-.27,.03) circle (.135);
\draw (.27,.03) circle (.135);
\draw (-.19,.14) .. controls (-.55,.78) and (.55,.78) .. (.19,.14);
\end{tikzpicture}}}}}
\clearpage\begingroup\NoSourceHead\SourceLeaf{226}
\begin{center}
{\Large ELLIPTISCHE BAHNBESTIMMUNG.}\par\vspace{1.2em}
{\large\textls[140]{NACHLASS.}}\par\vspace{1.2em}
[I.]
\end{center}
\noindent{\small [S. 1]}
\begin{center}
BESTIMMUNG DER BAHNEN DER HIMMELSKÖRPER.
\end{center}
\noindent\rule{\linewidth}{.35pt}\par
\noindent{\small [Aus Handbuch 16, Bb, Den astronomischen Wissenschaften gewidmet, November 1801, S. 1—8.]}
\par\noindent\rule{\linewidth}{.35pt}
\begin{center}
[1.]\par\medskip
\textls[100]{Bedeutung der hier vorkommenden Zeichen.}
\end{center}
\noindent\begin{minipage}[t]{.48\linewidth}
\small\setlength{\tabcolsep}{0pt}\renewcommand{\arraystretch}{1.2}
\begin{tabular}{@{}l@{\quad}p{.77\linewidth}@{}}
$\alpha$ .\,.\,. & Geocentrische Länge des Weltkörpers\\
$\delta$ .\,.\,. & Geocentrische Breite\\
$\Theta$ .\,.\,. & Tangente der Breite\\
$\Delta$ .\,.\,. & Kurtirter Abstand von der Erde\\
$\lambda$ .\,.\,. & Heliocentrische Länge\\
$\beta$ .\,.\,. & Heliocentrische Breite\\
$\vartheta$ .\,.\,. & Tangente dieser Breite\\
$\rho$ .\,.\,. & Kurtirter Abstand von der Sonne\\
$r$ .\,.\,. & Wahrer Abstand von der Sonne
\end{tabular}
\end{minipage}\hfill\vrule width .35pt\hfill
\begin{minipage}[t]{.48\linewidth}
\small\setlength{\tabcolsep}{0pt}\renewcommand{\arraystretch}{1.2}
\begin{tabular}{@{}l@{\quad}p{.77\linewidth}@{}}
$t$ .\,.\,. & Zeit der Beobachtung\\
$v$ .\,.\,. & Länge in der Bahn\\
$M$ .\,.\,. & Wahre Anomalie [\textsuperscript{*})]\\
$E$ .\,.\,. & Eccentrische Anomalie [\textsuperscript{*})]\\
$N$ .\,.\,. & Mittlere Anomalie [\textsuperscript{*})]\\
$u$ .\,.\,. & Mittlere Länge\\
$R$ .\,.\,. & Abstand der Sonne von der Erde\\
$V$ .\,.\,. & Wahre Länge der Sonne\\
$U$ .\,.\,. & Mittlere Länge der Sonne.
\end{tabular}
\end{minipage}
\par\begingroup\renewcommand{\thefootnote}{}\footnotetext{[\textsuperscript{*}) Vom Aphelium gezählt.]}\endgroup
\clearpage
\HistoricalHead{222}{L}{THEORETISCHE ASTRONOMIE. NACHLASS.}\SourceLeaf{227}
Dieselben Zeichen mit Linien haben ähnliche Bedeutungen für andere Beobachtungen.

\medskip\noindent\begin{minipage}[c]{.54\linewidth}
\setlength{\tabcolsep}{0pt}\renewcommand{\arraystretch}{1.2}
\begin{tabular}{@{}l@{\quad}p{.78\linewidth}@{}}
$\AstNode$ .\,.\,. & Länge des aufsteigenden Knoten\\
$i$ .\,.\,. & Neigung der Bahn\\
$g$ .\,.\,. & Tangente derselben\\
$\omega$ .\,.\,. & Ort der Sonnenferne\\
$e$ .\,.\,. & Eccentricität\\
$\varphi$ .\,.\,. & Bogen dessen Sinus $=e$\\
$a$ .\,.\,. & Halbe grosse Axe\\
$b$ .\,.\,. & Halbe kleine Axe\\
$p$ .\,.\,. & Halber Parameter
\end{tabular}
\end{minipage}\hfill
\begin{minipage}[c]{.42\linewidth}
\setlength{\parindent}{1.2em}
Anm. Die Längen sind sämmtlich siderisch, oder von einem festen, übrigens willkürlichen Punkte des Himmels an gezählt.
\end{minipage}
\clearpage\endgroup
